The perihelion distance is $r_p=a(1-e)$ and the aphelion distance is $r_a=a(1+e)$. \hfill(2 points)
The law of conservation of energy gives
\[E=-\frac{GMm}{r_p}+\frac{mv_p^2}{2}=-\frac{GMm}{r_a}+\frac{mv_a^2}{2}.\]
\hfill(2 points)
From conservation of angular momentum,
\[mv_ar_a=mv_pr_p,\qquad v_a=v_p\frac{r_p}{r_a}.\]
\hfill(1 point)
Substitute velocities and rearrange the first equation. An alternative expression making $v_p$ the subject is also correct.
\[v_a^2\left(1-\frac{r_a^2}{r_p^2}\right)=2GM\left(\frac1{r_a}-\frac1{r_p}\right).\]
\hfill(3 points)
Substitute into the equation for the total mechanical energy,
\[E=-\frac{GMm}{r_a}+\frac{GMmr_p^2}{r_p^2-r_a^2}\left(\frac{r_p-r_a}{r_ar_p}\right).\]
After simplifying or any equivalent formulation,
\[E=-\frac{GMm}{2a}\left(\frac2{1+e}-\frac{1-e}{1+e}\right).\]
\hfill(3 points)
Finally, $E=-GMm/(2a)$. \hfill(2 points)
Substituting $r_a$ and $r_p$ in the energy equations gives $v_a/v_p=(1-e)/(1+e)$. One can also use angular momentum directly. \hfill(2 points)
% editorial: The official key gives the reciprocal of the ratio requested in the statement. Its displayed relation is retained.
